\label{part:geometry}

\textbf{（表示几何：结构基底与联合状态）}

\bluetext{状态变化所依赖的组织结构}

前一篇讨论目标、历史和外部输入如何共同约束智能过程。本篇暂时固定任务与结构版本，研究状态更新所依赖的表示组织：我们把什么视为一个对象，关系如何定义可达路径，属性怎样归属于对象，不同坐标中的描述又怎样保持任务读出的一致性。固定结构是一项分析条件，不是说实际智能的结构永远不变。只要对象数量、关系类型或表示接口发生改变，我们就需要记录结构事件，并重新检查原模型的空间、维度和假设。

我们从语义单元开始。有限资源要求选择采样率、编码精度和任务相关粒度，却不要求把概念当成不可分割的自然粒子。形与质分别描述关系组织和属性内容，联合表示通过身份、绑定与读出接口建立。张量是一种可选择的数学实现，不是所有智能系统的唯一本体；同样，互斥分类来自任务契约，多关系组合来自类型系统，不借助量子统计来证明。我们关心的是从输入索引到对象记录的形成过程，以及错误归属如何被识别和纠正。

随后我们研究结构基底。图、超图和单纯复形可以表达不同阶数的关系，但每种关系都需要明确的语义和来源。图中的可达性不自动等同于因果关系，图上的距离也未必符合对称性或三角不等式。价值与经历可以通过权重、偏序、目标和证据影响传播，却不能只凭一种几何意象就被当作自然作用力。对固定结构得到的性质需要在更新时重新审查；删除节点、修改关系或合并身份之后，旧路径、缓存和引用是否有效，都属于模型的一部分。

本篇还将分析联合状态与表示变换。局部结构网络保持对象、部件及事件角色，全局分布表示支持共享模式与跨情境泛化；二者通过任务相关身份、绑定、耦合和读出联合，而不是互相替代。我们分别讨论坐标变化、活动传播、属性修订和结构编辑，避免把不同变化统称为同一种运动。一个量的保持必须来自明确定义与保持条件；若输入、目标、度量或结构随时间变化，它们应进入完整导数或离散更新关系。内部读出保持、事实正确和任务成功也需要分别验证。

为方便后续论证，我们规定连续时间变量按声明的实际时间尺度计量，离散下标表示采样或事务步骤；任务进展与记忆内化次序可以另设索引，但不能未经定义就当成第二种物理时间。属性向量在数值模型中按参考尺度归一化，原始测量值及单位保留在有类型字段里。活动、概率、目标代价、资源预算和实际能耗分别记录。AgentOS 中自我相关状态沿用 $\Psi$，一般状态场采用另外的符号，防止把工程部件与任意复数向量混同。

本篇的静态分析最终服务于动态模型：只有对象身份、属性归属、关系接口和有效域足够明确，后续传播、控制、记忆沉淀及结构适应才有可检查的作用对象。我们允许使用几何语言压缩论证，也允许在相应领域解释真实物理对象，但任何类比都不能代替数学条件或经验检验。HSF-HD在这里构建一般智能的状态模型族，AgentOS则提供可执行的工程实例；两者的层级保持不变。
